\subsection{معرفی و انگیزه}
این ماژول در چارچوب پژوهش حاضر پیاده‌سازی و ارزیابی شده و فاقد مقاله مستقل جداگانه در ادبیات تحقیق است.

هرس کردن وزن‌ها \lr{(Weight Pruning)} به عنوان یکی از موثرترین تکنیک‌های سبک‌سازی شبکه‌های عصبی شناخته می‌شود؛ با این حال، حذف ناگهانی یا مقداری خالص وزن‌ها \lr{(Magnitude Pruning)} معمولاً به از دست رفتن اتصالات حیاتی و زوال ظرفیت زبانی مدل می‌انجامد. علاوه بر این، تنک‌سازی نامنظم \lr{(Unstructured Sparsity)} به دلیل عدم تطابق با خطوط لوله حافظه در شتاب‌دهنده‌های گرافیکی، شتاب محاسباتی ملموسی در استنتاج ایجاد نمی‌کند.

ماژول فشرده‌سازی تخصصی \lr{ADMMPruner} یک چارچوب بهینه‌سازی دقیق مبتنی بر الگوریتم جهت‌های متناوب لاگرانژی افزوده \lr{(ADMM)} را به همراه تقطیر بازسازی لایه‌ای \lr{(Layer-Wise Reconstruction Distillation)}، تخمین‌گر گرادیان مستقیم \lr{(Straight-Through Estimator - STE)} و الگوی تنک‌سازی ساختاریافته ۲:۴ معرفی می‌نماید. در این ساختار، متغیر تنک در تعاملی پویا با وزن‌های پیوسته به سمت آرایش ماتریسی هرس‌شده با کمترین اتلاف خروجی هدایت می‌شود و با به‌روزرسانی پویای ضریب جریمه $\rho$، پایداری همگرایی تضمین می‌گردد.
